\begin{abstract}


Semantic segmentation usually suffers from a long-tail data distribution.  % , which could cause biased predictions.
%
Due to the imbalanced number of samples across categories, the features of those tail classes may get squeezed into a narrow area in the feature space.  % , leaving the decision boundaries hard to learn.
%
Towards a balanced feature distribution, we introduce category-wise variation into the network predictions in the training phase such that an instance is no longer projected to a feature point, but a small region instead.
%
Such a perturbation is highly dependent on the category scale, which appears as assigning smaller variation to head classes and larger variation to tail classes.
%
In this way, we manage to close the gap between the feature areas of different categories, resulting in a more balanced representation.  % and hence alleviate the learning difficulty.
%
It is noteworthy that the introduced variation is discarded at the inference stage to facilitate a confident prediction.
%
Although with an embarrassingly simple implementation, our method manifests itself in strong generalizability to various datasets and task settings.
%
Extensive experiments suggest that our plug-in design lends itself well to a range of state-of-the-art approaches and boosts the performance on top of them.%
%
\footnote{
Code: \url{https://github.com/grantword8/BLV}.

~$^\dag$This work was done during the internship at SenseTime Research.

~$^\ddagger$Rui Zhao is also with Qing Yuan Research Institute, Shanghai Jiao Tong University.
}


\iffalse
%
Due to the long-tailed distribution of training set, semantic segmentation model is severely biased towards the dominant category. 
%
We observe that vanilla training on long-tailed data with cross-entropy loss makes the instance-rich head classes severely squeeze the spatial distribution of the tail classes, which leads to insufficient training of tail class pixels. 
%
It is unfavorable for training on balanced data, but can be utilized to adjust the validity of the samples in long-tailed data, thereby solving the distorted embedding
space of long-tailed traning data. 
%
To this end, this paper proposes the a novel loss called "Augment Your Logits"(\method) to augment the final logits output of the training model.
%
We set relatively large the amplitude of perturbation to the tail class logits and small that of the head class.
%
The large cloud size can reduce the softmax saturation and thereby making tail class samples more active as well as enlarging the embedding space. 
%
Extensive experiments on benchmark datasets and three main stream semantic segmentation tasks validate the superior performance of the proposed method. 
%
Source code will be made public available.
%
\fi


\end{abstract}
